% Options for packages loaded elsewhere
\PassOptionsToPackage{unicode}{hyperref}
\PassOptionsToPackage{hyphens}{url}
\documentclass[
]{article}
\usepackage{xcolor}
\usepackage[margin=1in]{geometry}
\usepackage{amsmath,amssymb}
\setcounter{secnumdepth}{-\maxdimen} % remove section numbering
\usepackage{iftex}
\ifPDFTeX
  \usepackage[T1]{fontenc}
  \usepackage[utf8]{inputenc}
  \usepackage{textcomp} % provide euro and other symbols
\else % if luatex or xetex
  \usepackage{unicode-math} % this also loads fontspec
  \defaultfontfeatures{Scale=MatchLowercase}
  \defaultfontfeatures[\rmfamily]{Ligatures=TeX,Scale=1}
\fi
\usepackage{lmodern}
\ifPDFTeX\else
  % xetex/luatex font selection
\fi
% Use upquote if available, for straight quotes in verbatim environments
\IfFileExists{upquote.sty}{\usepackage{upquote}}{}
\IfFileExists{microtype.sty}{% use microtype if available
  \usepackage[]{microtype}
  \UseMicrotypeSet[protrusion]{basicmath} % disable protrusion for tt fonts
}{}
\makeatletter
\@ifundefined{KOMAClassName}{% if non-KOMA class
  \IfFileExists{parskip.sty}{%
    \usepackage{parskip}
  }{% else
    \setlength{\parindent}{0pt}
    \setlength{\parskip}{6pt plus 2pt minus 1pt}}
}{% if KOMA class
  \KOMAoptions{parskip=half}}
\makeatother
\usepackage{color}
\usepackage{fancyvrb}
\newcommand{\VerbBar}{|}
\newcommand{\VERB}{\Verb[commandchars=\\\{\}]}
\DefineVerbatimEnvironment{Highlighting}{Verbatim}{commandchars=\\\{\}}
% Add ',fontsize=\small' for more characters per line
\usepackage{framed}
\definecolor{shadecolor}{RGB}{248,248,248}
\newenvironment{Shaded}{\begin{snugshade}}{\end{snugshade}}
\newcommand{\AlertTok}[1]{\textcolor[rgb]{0.94,0.16,0.16}{#1}}
\newcommand{\AnnotationTok}[1]{\textcolor[rgb]{0.56,0.35,0.01}{\textbf{\textit{#1}}}}
\newcommand{\AttributeTok}[1]{\textcolor[rgb]{0.13,0.29,0.53}{#1}}
\newcommand{\BaseNTok}[1]{\textcolor[rgb]{0.00,0.00,0.81}{#1}}
\newcommand{\BuiltInTok}[1]{#1}
\newcommand{\CharTok}[1]{\textcolor[rgb]{0.31,0.60,0.02}{#1}}
\newcommand{\CommentTok}[1]{\textcolor[rgb]{0.56,0.35,0.01}{\textit{#1}}}
\newcommand{\CommentVarTok}[1]{\textcolor[rgb]{0.56,0.35,0.01}{\textbf{\textit{#1}}}}
\newcommand{\ConstantTok}[1]{\textcolor[rgb]{0.56,0.35,0.01}{#1}}
\newcommand{\ControlFlowTok}[1]{\textcolor[rgb]{0.13,0.29,0.53}{\textbf{#1}}}
\newcommand{\DataTypeTok}[1]{\textcolor[rgb]{0.13,0.29,0.53}{#1}}
\newcommand{\DecValTok}[1]{\textcolor[rgb]{0.00,0.00,0.81}{#1}}
\newcommand{\DocumentationTok}[1]{\textcolor[rgb]{0.56,0.35,0.01}{\textbf{\textit{#1}}}}
\newcommand{\ErrorTok}[1]{\textcolor[rgb]{0.64,0.00,0.00}{\textbf{#1}}}
\newcommand{\ExtensionTok}[1]{#1}
\newcommand{\FloatTok}[1]{\textcolor[rgb]{0.00,0.00,0.81}{#1}}
\newcommand{\FunctionTok}[1]{\textcolor[rgb]{0.13,0.29,0.53}{\textbf{#1}}}
\newcommand{\ImportTok}[1]{#1}
\newcommand{\InformationTok}[1]{\textcolor[rgb]{0.56,0.35,0.01}{\textbf{\textit{#1}}}}
\newcommand{\KeywordTok}[1]{\textcolor[rgb]{0.13,0.29,0.53}{\textbf{#1}}}
\newcommand{\NormalTok}[1]{#1}
\newcommand{\OperatorTok}[1]{\textcolor[rgb]{0.81,0.36,0.00}{\textbf{#1}}}
\newcommand{\OtherTok}[1]{\textcolor[rgb]{0.56,0.35,0.01}{#1}}
\newcommand{\PreprocessorTok}[1]{\textcolor[rgb]{0.56,0.35,0.01}{\textit{#1}}}
\newcommand{\RegionMarkerTok}[1]{#1}
\newcommand{\SpecialCharTok}[1]{\textcolor[rgb]{0.81,0.36,0.00}{\textbf{#1}}}
\newcommand{\SpecialStringTok}[1]{\textcolor[rgb]{0.31,0.60,0.02}{#1}}
\newcommand{\StringTok}[1]{\textcolor[rgb]{0.31,0.60,0.02}{#1}}
\newcommand{\VariableTok}[1]{\textcolor[rgb]{0.00,0.00,0.00}{#1}}
\newcommand{\VerbatimStringTok}[1]{\textcolor[rgb]{0.31,0.60,0.02}{#1}}
\newcommand{\WarningTok}[1]{\textcolor[rgb]{0.56,0.35,0.01}{\textbf{\textit{#1}}}}
\usepackage{graphicx}
\makeatletter
\newsavebox\pandoc@box
\newcommand*\pandocbounded[1]{% scales image to fit in text height/width
  \sbox\pandoc@box{#1}%
  \Gscale@div\@tempa{\textheight}{\dimexpr\ht\pandoc@box+\dp\pandoc@box\relax}%
  \Gscale@div\@tempb{\linewidth}{\wd\pandoc@box}%
  \ifdim\@tempb\p@<\@tempa\p@\let\@tempa\@tempb\fi% select the smaller of both
  \ifdim\@tempa\p@<\p@\scalebox{\@tempa}{\usebox\pandoc@box}%
  \else\usebox{\pandoc@box}%
  \fi%
}
% Set default figure placement to htbp
\def\fps@figure{htbp}
\makeatother
\setlength{\emergencystretch}{3em} % prevent overfull lines
\providecommand{\tightlist}{%
  \setlength{\itemsep}{0pt}\setlength{\parskip}{0pt}}
\usepackage{bookmark}
\IfFileExists{xurl.sty}{\usepackage{xurl}}{} % add URL line breaks if available
\urlstyle{same}
\hypersetup{
  pdftitle={The Normal-Normal Model},
  hidelinks,
  pdfcreator={LaTeX via pandoc}}

\title{The Normal-Normal Model}
\author{}
\date{\vspace{-2.5em}2026-04-17}

\begin{document}
\maketitle

\subsection{Introduction}\label{introduction}

The Normal-Normal model is the simplest example of conjugate Bayesian
updating for a continuous parameter: a Normal prior on the mean, a
Normal likelihood with known variance, and a Normal posterior obtained
by a transparent precision-weighted average. It is the classical
pedagogical example, and despite its simplicity it sits at the heart of
every hierarchical model that contains a Normal level --- random-effects
meta-analysis, partial pooling in mixed models, and the mean structure
of Gaussian processes all rely on the same closed-form update.

\subsection{Prerequisites}\label{prerequisites}

A working understanding of the Normal distribution, the concept of
conjugacy, and the difference between variance and precision (inverse
variance). Basic algebra is enough to follow the derivation.

\subsection{Theory}\label{theory}

With prior \(\mu \sim \mathcal N(\mu_0, \tau_0^2)\) and likelihood
\(y_i \sim \mathcal N(\mu, \sigma^2)\) for \(i = 1, \dots, n\) with
known \(\sigma\), the posterior is

\[\mu \mid y \sim \mathcal N(\mu_n, \tau_n^2),\]

where the posterior precision is the sum of the prior precision and the
data precision,

\[\tau_n^{-2} = \tau_0^{-2} + n / \sigma^2,\]

and the posterior mean is the precision-weighted average

\[\mu_n = \tau_n^2 \left(\tau_0^{-2} \mu_0 + n \bar y / \sigma^2\right).\]

The posterior mean lies between the prior mean and the sample mean, with
weight that depends on relative precisions. The precision-additive
structure makes the prior's ``pseudo-sample-size'' interpretable:
\(\tau_0^2\) corresponds to a prior sample of size
\(\sigma^2 / \tau_0^2\).

\subsection{Assumptions}\label{assumptions}

Known \(\sigma^2\) --- so the only unknown is the mean --- and a Normal
prior. If the variance is also unknown, the conjugate setup uses a
Normal-Inverse-Gamma prior and the posterior on \(\mu\) is a \(t\)
distribution.

\subsection{R Implementation}\label{r-implementation}

\begin{Shaded}
\begin{Highlighting}[]
\CommentTok{\# Prior: mu \textasciitilde{} N(100, 10\^{}2); data: xbar = 105, sigma = 15, n = 20}
\NormalTok{mu0 }\OtherTok{\textless{}{-}} \DecValTok{100}\NormalTok{; tau0 }\OtherTok{\textless{}{-}} \DecValTok{10}
\NormalTok{xbar }\OtherTok{\textless{}{-}} \DecValTok{105}\NormalTok{; sigma }\OtherTok{\textless{}{-}} \DecValTok{15}\NormalTok{; n }\OtherTok{\textless{}{-}} \DecValTok{20}

\NormalTok{prec\_post }\OtherTok{\textless{}{-}} \DecValTok{1}\SpecialCharTok{/}\NormalTok{tau0}\SpecialCharTok{\^{}}\DecValTok{2} \SpecialCharTok{+}\NormalTok{ n}\SpecialCharTok{/}\NormalTok{sigma}\SpecialCharTok{\^{}}\DecValTok{2}
\NormalTok{tau\_post  }\OtherTok{\textless{}{-}} \DecValTok{1}\SpecialCharTok{/}\FunctionTok{sqrt}\NormalTok{(prec\_post)}
\NormalTok{mu\_post   }\OtherTok{\textless{}{-}}\NormalTok{ tau\_post}\SpecialCharTok{\^{}}\DecValTok{2} \SpecialCharTok{*}\NormalTok{ (mu0}\SpecialCharTok{/}\NormalTok{tau0}\SpecialCharTok{\^{}}\DecValTok{2} \SpecialCharTok{+}\NormalTok{ n}\SpecialCharTok{*}\NormalTok{xbar}\SpecialCharTok{/}\NormalTok{sigma}\SpecialCharTok{\^{}}\DecValTok{2}\NormalTok{)}

\FunctionTok{c}\NormalTok{(}\AttributeTok{posterior\_mean =}\NormalTok{ mu\_post,}
  \AttributeTok{posterior\_sd   =}\NormalTok{ tau\_post,}
  \AttributeTok{ci95 =}\NormalTok{ mu\_post }\SpecialCharTok{+} \FunctionTok{c}\NormalTok{(}\SpecialCharTok{{-}}\DecValTok{1}\NormalTok{, }\DecValTok{1}\NormalTok{) }\SpecialCharTok{*} \FloatTok{1.96} \SpecialCharTok{*}\NormalTok{ tau\_post)}
\end{Highlighting}
\end{Shaded}

\subsection{Output \& Results}\label{output-results}

Posterior mean approximately 104.2 with posterior standard deviation 2.4
and 95 \% credible interval roughly \((99.6, 108.9)\). The shrinkage
from the sample mean of 105 toward the prior mean of 100 is small
because the data precision (\(n / \sigma^2 \approx 0.089\)) is much
larger than the prior precision (\(1 / \tau_0^2 = 0.01\)).

\subsection{Interpretation}\label{interpretation}

A reporting sentence: ``Combining a \(\mathcal N(100, 10^2)\) prior with
20 observations of sample mean 105 (known \(\sigma = 15\)) yielded a
posterior mean of 104.2 (95 \% credible interval 99.6 to 108.9); the
prior contributed the equivalent of approximately 2.25 additional
observations.'' The ``equivalent observations'' framing is one of the
most intuitive ways to communicate prior strength to non-Bayesians.

\subsection{Practical Tips}\label{practical-tips}

\begin{itemize}
\tightlist
\item
  As \(n \to \infty\), the posterior mean tends to the sample mean and
  the posterior variance tends to \(\sigma^2 / n\); the prior's
  influence vanishes asymptotically but can dominate in small samples.
\item
  When \(\sigma\) is unknown, switch to the Normal--Inverse-Gamma
  conjugate pair; the marginal posterior on \(\mu\) then has \(t\)
  tails.
\item
  The precision-form update generalises directly to multivariate Normals
  with covariance matrices in place of variances; everything that was
  scalar becomes matrix.
\item
  Use this model as a sanity check for software: any general-purpose
  Bayesian engine (Stan, JAGS, \texttt{brms}) should reproduce the
  closed-form posterior to within Monte Carlo error.
\item
  Hierarchical mean structures are built from this update at every
  level, which is why understanding it deeply pays dividends in mixed
  models and meta-analysis.
\end{itemize}

\subsection{R Packages Used}\label{r-packages-used}

Base R is sufficient because the update is closed-form (\texttt{dnorm},
\texttt{pnorm}, \texttt{qnorm}); for didactic plotting compare prior,
likelihood, and posterior with \texttt{ggplot2}, and for general
Bayesian alternatives any of \texttt{brms}, \texttt{rstanarm}, or
\texttt{BayesFactor} reproduce the same answer.

\subsection{For Reviewers}\label{for-reviewers}

\textbf{What to look for in a paper using this method.}

\begin{itemize}
\tightlist
\item
  \textbf{Common misapplications.}
\item
  \textbf{Diagnostics that should be reported but often aren't.}
\item
  \textbf{Red flags in tables and figures.}
\item
  \textbf{What to verify.}
\item
  \textbf{What an adequate Methods paragraph must contain.}
\end{itemize}

\subsection{See also --- labs in this
chapter}\label{see-also-labs-in-this-chapter}

\begin{itemize}
\tightlist
\item
  \href{labs/lab-bayesian-thinking-control-vs-decision-error.Rmd}{Bayesian
  thinking; α control vs decision error}
\item
  \href{labs/lab-brms-stan-loo-hierarchical-models.Rmd}{brms / Stan,
  LOO, hierarchical models}
\end{itemize}

\end{document}
